\chapter*{ВВЕДЕНИЕ}
\addcontentsline{toc}{chapter}{ВВЕДЕНИЕ}

Развитие систем охранного видеонаблюдения сопровождается постоянным ростом числа камер, объема поступающих видеоданных и требований к скорости реагирования на потенциальные инциденты. Современные распределенные объекты охраны (промышленные площадки, транспортные узлы, периметры закрытых территорий, объекты критической информационной инфраструктуры) ежедневно генерируют значительный объем видеоматериала, существенная часть которого требует оперативной интерпретации в режиме реального времени. В таких условиях ручной мониторинг видеопотоков становится организационно ограниченным, а практическая ценность платформы определяется сочетанием автоматического обнаружения, сопровождения объектов и надежной событийной доставки~\cite{redmon2016yolo,wojke2017deepsort,transactional_outbox}.

Вместе с тем требования нормативных и отраслевых документов, регулирующих проектирование, испытания и эксплуатацию систем видеонаблюдения, формируют дополнительные ограничения на полноту фиксации событий, время их доставки во внешние системы и устойчивость к отказам отдельных компонентов~\cite{iec_62676_4_2025,iec_62676_6_2026}. Это закономерно повышает потребность в платформах видеоаналитики реального времени, способных автоматически обнаруживать объекты, сопровождать их в последовательности кадров, интерпретировать наблюдаемую ситуацию и надежно передавать сведения о событиях внешним потребителям. Особую значимость приобретают сценарии, связанные с обнаружением малозаметных и динамичных объектов, в том числе беспилотных летательных аппаратов, которые на ряде объектов охраны рассматриваются как самостоятельный класс угроз~\cite{chen2026uavdb}.

\textbf{Актуальность темы исследования} определяется как научными, так и практическими факторами. С научной точки зрения существенный интерес представляет комплексное согласование методов обнаружения объектов, их сопровождения и событийного анализа с архитектурными механизмами, обеспечивающими работу системы в условиях ограниченной задержки и неустойчивых внешних интеграций. Для целей настоящей работы недостаточно рассматривать только качество детектора или только надежность программной инфраструктуры: значимым является их совместное проектирование и экспериментальная проверка в едином контуре. С практической точки зрения востребованы решения, позволяющие сократить нагрузку на операторов служб безопасности, уменьшить время реакции на инциденты, повысить надежность мониторинга, обеспечить совместимость с действующими системами оповещения и управления и снизить риск потери событий при сбоях транспортного слоя~\cite{transactional_outbox,prometheus_docs}.

\textbf{Цель работы} состоит в разработке и экспериментальной проверке на исследовательско-экспериментальном стенде платформы видеоаналитики реального времени для мониторинга охраняемых зон и событийного анализа, обеспечивающей обнаружение и сопровождение объектов, формирование интерпретируемых событий и надежную доставку сведений о событиях внешним потребителям.

Для достижения поставленной цели необходимо решить следующие задачи:
\begin{enumerate}
    \item проанализировать предметную область охранного видеонаблюдения, существующие подходы к построению платформ видеоаналитики и сформировать функциональные и нефункциональные требования к разрабатываемому решению;
    \item обосновать выбор методов обнаружения объектов, сопровождения объектов и правилового событийного анализа для целевых сценариев видеонаблюдения, включая демонстрационный сценарий обнаружения беспилотных летательных аппаратов;
    \item разработать модель, архитектуру и алгоритмический конвейер платформы, объединяющие прием видеопотока, обнаружение, сопровождение, формирование событий и асинхронную доставку сведений о событиях внешним потребителям;
    \item реализовать программную платформу с компонентами управления, аналитической обработки, интеграционной доставки, разграничения доступа, наблюдаемости и воспроизводимого развертывания;
    \item провести стендовую апробацию платформы по интеграционным, отказным, ролевым и демонстрационным видеосценариям с фиксацией границ доказанности полученных результатов;
    \item определить ограничения текущей реализации, условия ее применимости и направления дальнейшего развития для перехода от исследовательско-экспериментального стенда к промышленному применению.
\end{enumerate}

\textbf{Объектом исследования} являются процессы автоматизированного мониторинга видеопотоков в системах охранного видеонаблюдения.

\textbf{Предмет исследования} составляют методы, модели и архитектурные решения платформ видеоаналитики реального времени, обеспечивающие обнаружение и сопровождение объектов, формирование событий и надежную доставку сведений о событиях при ограничениях по задержке обработки, отказоустойчивости и безопасности доступа.

\textbf{Теоретическую и методическую основу исследования} составляют:
\begin{itemize}
    \item методы компьютерного зрения и машинного обучения для задач обнаружения и сопровождения объектов в видеопотоках~\cite{redmon2016yolo,wojke2017deepsort,kalman1960};
    \item методы цифровой обработки видео, включая декодирование потоков и работу с временными метками~\cite{ffmpeg_docs,pyav_docs};
    \item методы оценки качества алгоритмов обнаружения и сопровождения, в том числе на основе методик MOT-Challenge и метрик MOTMetrics~\cite{bernardin2008motmetrics,motchallenge};
    \item методы проектирования распределенных программных систем и надежной асинхронной доставки сообщений, в том числе шаблон транзакционной таблицы исходящих сообщений и политики повторной доставки~\cite{richardson_microservices,transactional_outbox,rabbitmq_dlq};
    \item методы экспериментальной оценки характеристик систем реального времени, включая измерение сквозной задержки, контроль устойчивости при отказах и сбор операционных метрик~\cite{prometheus_docs,grafana_docs,iec_62676_6_2026};
    \item подходы к разграничению доступа и аудиту действий пользователей в информационных системах~\cite{nist_sp800_53r5};
    \item требования и подходы к оформлению результатов научно-технических исследований, в том числе ГОСТ~7.32-2017~\cite{gost_732_2017}, ГОСТ~Р~7.0.5-2008~\cite{gost_r_705_2008} и сопутствующие нормативные документы~\cite{gost_71_2003,gost_780_2000,gost_782_2001}.
\end{itemize}

\textbf{Научная новизна работы} заключается в следующем:
\begin{enumerate}
    \item предложен и стендово проверен архитектурно-алгоритмический подход к построению платформы видеоаналитики реального времени для задач охранного видеонаблюдения, обеспечивающий согласованное решение задач обнаружения, сопровождения, событийной интерпретации и надежной доставки сведений о событиях в едином программном контуре;
    \item уточнена модель сквозной задержки от объекта к событию (\(L_{\texttt{object} \rightarrow \texttt{event}}\)) с разделением на этапы декодирования, инференса, сопровождения, применения правил и асинхронной публикации, что позволяет локализовать узкие места без изменения внешних контрактов системы;
    \item сформулированы и реализованы инженерные принципы интеграции аналитического ядра с внешними потребителями событий через шаблон транзакционной таблицы исходящих сообщений, брокер сообщений и контролируемую политику повторов, обеспечивающие сохранность доменных событий в условиях кратковременных отказов внешних компонентов;
    \item разработан набор эксплуатационных проверок (интеграционные сценарии доставки, отказные сценарии RabbitMQ и HTTP-получателя, сценарии RBAC/JWT), пригодных для регулярного контроля архитектурных свойств платформы между релизами.
\end{enumerate}

\textbf{Положения, выносимые на защиту:}
\begin{enumerate}
    \item архитектурно-алгоритмический подход к построению платформы видеоаналитики реального времени, разделяющий контуры управления и обработки данных и обеспечивающий совместное решение задач обнаружения, сопровождения и событийного анализа объектов;
    \item модель сквозной задержки и методика ее декомпозиции, позволяющая обоснованно выбирать направления оптимизации горячего пути обработки без нарушения внешних контрактов системы;
    \item инженерное решение надежной асинхронной доставки доменных событий на основе шаблона транзакционной таблицы исходящих сообщений и контролируемых повторов с маршрутизацией в очередь недоставленных сообщений;
    \item методика экспериментальной приемки платформы по совокупности интеграционных, отказных и ролевых сценариев проверки безопасности с фиксацией ключевых эксплуатационных метрик.
\end{enumerate}

\textbf{Практическая значимость работы} состоит в возможности применения разработанной платформы для повышения оперативности выявления инцидентов, снижения нагрузки на операторов служб безопасности и интеграции интеллектуальной видеоаналитики с действующими системами оповещения и управления безопасностью. Результаты исследования могут быть использованы при создании и модернизации систем мониторинга охраняемых зон на объектах с повышенными требованиями к непрерывности наблюдения и скорости реакции, а также при подготовке технических заданий и регламентов эксплуатации платформ видеоаналитики, в том числе в части требований к надежности доставки событий и наблюдаемости.

\textbf{Методология и методы исследования.} В работе применен системный подход к проектированию программно-алгоритмических систем реального времени. Использовались методы анализа требований, проектирования архитектуры распределенных систем, экспериментального тестирования и сбора эксплуатационных метрик. Программная реализация выполнена с применением технологий FastAPI, PostgreSQL, RabbitMQ, Docker Compose, Prometheus и Grafana~\cite{fastapi_docs,sqlalchemy_docs,rabbitmq_docs,docker_docs,prometheus_docs,grafana_docs}. Алгоритмическая часть основана на одноэтапных детекторах семейства YOLO и модулях сопровождения семейства SORT/DeepSORT с дополнительной обработкой результатов в правиловом событийном движке~\cite{redmon2016yolo,bewley2016sort,wojke2017deepsort}.

\textbf{Апробация результатов работы.} Работоспособность и устойчивость платформы подтверждены автоматизированным приемочным прогоном интеграционного контура на CPU-профиле, регрессионными тестами API и фоновых обработчиков, а также демонстрационным сценарием обнаружения беспилотных летательных аппаратов. В рамках апробации зафиксированы:
\begin{itemize}
    \item прохождение синтетически зафиксированных событий \texttt{zone\_enter}, \texttt{line\_cross} и \texttt{dwell\_time} через транзакционную таблицу исходящих сообщений, брокер сообщений и внешний HTTP-получатель;
    \item демонстрационный путь обработки БПЛА-видео от обученной модели до траекторий, событий и операторского интерфейса;
    \item устойчивость интеграционного контура при кратковременной недоступности RabbitMQ и HTTP-получателя, отсутствие потерь событий и завершение доставки после восстановления внешних компонентов;
    \item соблюдение ролевой модели доступа \texttt{admin/operator/viewer} в сценариях управления и чтения данных;
    \item инструментальная готовность к сбору операционных метрик средствами Prometheus и их визуализации в Grafana; статистически устойчивая p95-оценка полного пути от объекта к событию вынесена в последующий длительный стендовый прогон.
\end{itemize}

Основной проверочный артефакт апробации -- сохраненный JSON-файл приемочного прогона от 13.05.2026; он используется в работе как источник статусов, счетчиков доставки и времен восстановления, а не как иллюстративное приложение. Демонстрационные материалы БПЛА-профиля используются отдельно для подтверждения связности видеоконтура. Сценарий исчерпания повторов и перехода в очередь недоставленных сообщений подтвержден отдельной проверкой логики фонового обработчика и не смешивается с результатами RUN-2026-05-13-A1.

Таким образом, в работе не отождествляются три уровня проверки: приемочный прогон интеграционного контура, демонстрационный видеосценарий и промышленная многокамерная оценка задержки. Количественные выводы формулируются только в пределах фактически сохраненных артефактов.

Подготовлен экспериментальный протокол, предназначенный для повторения CPU-проверок, фиксации отказных сценариев и подготовки публикации по итогам исследования. Программная реализация платформы оформлена в виде набора микросервисов с воспроизводимым стендом развертывания.

\textbf{Структура и объем работы.} Работа состоит из введения, четырех глав, заключения, списка использованных источников и приложений. Во введении обоснована актуальность темы, сформулированы цель, задачи, объект, предмет и научная новизна работы, изложена методологическая основа и приведены сведения об апробации. В первой главе выполнен анализ предметной области, рассмотрены существующие подходы к построению платформ охранной видеоаналитики и сформулированы требования к разрабатываемому решению. Вторая глава посвящена теоретическим предпосылкам, моделям и методологии: рассматриваются логика работы платформы, ее архитектура и используемый технологический стек, а также ключевые алгоритмы и расчетные соотношения. В третьей главе приведена программная реализация платформы, включая API, фоновые обработчики, пользовательский интерфейс (UI), контракты, листинги ключевых фрагментов кода и эксплуатационные регламенты. В четвертой главе представлены результаты апробации, таблицы приемочного прогона, демонстрационный БПЛА-профиль, графики обучения детектора и ограничения текущих измерений. В заключении систематизированы основные выводы и обозначены перспективы дальнейшего развития.
